% ---- pg-070 （书页 58） ----

注意有 $\bm{\Sigma}_{12} = \bm{\Sigma}_{21}'$, $\bm{X}^{(1)}$ 的协方差阵是 $\bm{\Sigma}_{11}$, $\bm{X}^{(2)}$ 的是 $\bm{\Sigma}_{22}$,且来自 $\bm{X}^{(1)}$ 和 $\bm{X}^{(2)}$ 元素的协方差阵是 $\bm{\Sigma}_{12}$(或 $\bm{\Sigma}_{21}$).

有时候用记号 $\operatorname{Cov}(\bm{X}^{(1)},\bm{X}^{(2)})$ 更为方便些,这里

\[
\operatorname{Cov}(\bm{X}^{(1)},\bm{X}^{(2)}) = \bm{\Sigma}_{12}
\]

是一个由 $\bm{X}^{(1)}$ 和 $\bm{X}^{(2)}$ 的分量之间的协方差构成的矩阵.

\mvsub{随机变量线性组合的均值向量和协方差阵}

如果单个随机变量,如 $X_1$,乘以一个常数 $c$,则有

\[
E(cX_1) = cE(X_1) = c\mu_1
\]

和

\[
\operatorname{Var}(cX_1) = E(cX_1 - c\mu_1)^2 = c^2\operatorname{Var}(X_1) = c^2\sigma_{11}
\]

如果 $X_2$ 是另一个随机变量且 $a, b$ 为常数,则利用期望的可加性,有

\begin{equation*}
\begin{aligned}
\operatorname{Cov}(aX_1,bX_2) &= E(aX_1 - a\mu_1)(bX_2 - b\mu_2) \\
&= abE(X_1 - \mu_1)(X_2 - \mu_2) \\
&= ab\operatorname{Cov}(X_1,X_2) = ab\sigma_{12}
\end{aligned}
\end{equation*}

最后,对于线性组合 $aX_1 + bX_2$,有

\phantomsection\label{eq:2-41}
\begin{equation}
\begin{aligned}
E(aX_1 + bX_2) &= aE(X_1) + bE(X_2) = a\mu_1 + b\mu_2 \\
\operatorname{Var}(aX_1 + bX_2) &= E[(aX_1 + bX_2) - (a\mu_1 + b\mu_2)]^2 \\
&= E[a(X_1 - \mu_1) + b(X_2 - \mu_2)]^2 \\
&= E\big[a^2(X_1 - \mu_1)^2 + b^2(X_2 - \mu_2)^2 + 2ab(X_1 - \mu_1)(X_2 - \mu_2)\big] \\
&= a^2\operatorname{Var}(X_1) + b^2\operatorname{Var}(X_2) + 2ab\operatorname{Cov}(X_1,X_2) \\
&= a^2\sigma_{11} + b^2\sigma_{22} + 2ab\sigma_{12}
\end{aligned} \tag{2-41}
\end{equation}

引入 $\bm{c}' = [a,b]$, $aX_1 + bX_2$ 可写为

\[
[a \quad b]\begin{bmatrix} X_1 \\ X_2 \end{bmatrix} = \bm{c}'\bm{X}
\]

类似的, $E(aX_1 + bX_2) = a\mu_1 + b\mu_2$ 可表示为

\[
[a \quad b]\begin{bmatrix} \mu_1 \\ \mu_2 \end{bmatrix} = \bm{c}'\bm{\mu}
\]

如果设

\[
\bm{\Sigma} = \begin{bmatrix} \sigma_{11} & \sigma_{12} \\ \sigma_{12} & \sigma_{22} \end{bmatrix}
\]

是 $\bm{X}$ 的方差-协方差矩阵,由于

\[
\bm{c}'\bm{\Sigma}\,\bm{c} = [a \quad b]\begin{bmatrix} \sigma_{11} & \sigma_{12} \\ \sigma_{12} & \sigma_{22} \end{bmatrix}\begin{bmatrix} a \\ b \end{bmatrix} = a^2\sigma_{11} + 2ab\sigma_{12} + b^2\sigma_{22}
\]

\hyperref[eq:2-41]{式(2-41)}变为

\phantomsection\label{eq:2-42}
\begin{equation}
\operatorname{Var}(aX_1 + bX_2) = \operatorname{Var}(\bm{c}'\bm{X}) = \bm{c}'\bm{\Sigma}\,\bm{c} \tag{2-42}
\end{equation}

上述结论可推广到 $p$ 个随机变量的线性组合：
